\documentclass[10pt,a4paper]{amsart}
\usepackage{amsmath,amssymb,amsthm,graphicx}
% For commutative diagrams we use Xy-pic
\usepackage[all]{xy}

% We use 2cell for 2-commutative diagrams.
\xyoption{2cell}
% Theorem environments.
%
\theoremstyle{plain}
\newtheorem{theorem}[subsection]{Theorem}
\newtheorem{proposition}[subsection]{Proposition}
\newtheorem{lemma}[subsection]{Lemma}

\theoremstyle{definition}
\newtheorem{definition}[subsection]{Definition}
\newtheorem{example}[subsection]{Example}
\newtheorem{exercise}[subsection]{Exercise}
\newtheorem{situation}[subsection]{Situation}

\theoremstyle{remark}
\newtheorem{remark}[subsection]{Remark}
\newtheorem{remarks}[subsection]{Remarks}

\numberwithin{equation}{subsection}

% Macros
%
\def\lim{\mathop{\mathrm{lim}}\nolimits}
\def\colim{\mathop{\mathrm{colim}}\nolimits}
\def\Spec{\mathop{\mathrm{Spec}}}
\def\Hom{\mathop{\mathrm{Hom}}\nolimits}
\def\Ext{\mathop{\mathrm{Ext}}\nolimits}
\def\SheafHom{\mathop{\mathcal{H}\!\mathit{om}}\nolimits}
\def\SheafExt{\mathop{\mathcal{E}\!\mathit{xt}}\nolimits}
\def\Sch{\mathit{Sch}}
\def\Mor{\mathop{\mathrm{Mor}}\nolimits}
\def\Ob{\mathop{\mathrm{Ob}}\nolimits}
\def\Sh{\mathop{\mathit{Sh}}\nolimits}
\def\NL{\mathop{N\!L}\nolimits}
\def\CH{\mathop{\mathrm{CH}}\nolimits}
\def\proetale{{pro\text{-}\acute{e}tale}}
\def\etale{{\acute{e}tale}}
\def\QCoh{\mathit{QCoh}}
\def\Ker{\mathop{\mathrm{Ker}}}
\def\Im{\mathop{\mathrm{Im}}}
\def\Coker{\mathop{\mathrm{Coker}}}
\def\Coim{\mathop{\mathrm{Coim}}}

% Boxtimes
%
\DeclareMathSymbol{\boxtimes}{\mathbin}{AMSa}{"02}

%
% Macros for moduli stacks/spaces
%
\def\QCohstack{\mathcal{QC}\!\mathit{oh}}
\def\Cohstack{\mathcal{C}\!\mathit{oh}}
\def\Spacesstack{\mathcal{S}\!\mathit{paces}}
\def\Quotfunctor{\mathrm{Quot}}
\def\Hilbfunctor{\mathrm{Hilb}}
\def\Curvesstack{\mathcal{C}\!\mathit{urves}}
\def\Polarizedstack{\mathcal{P}\!\mathit{olarized}}
\def\Complexesstack{\mathcal{C}\!\mathit{omplexes}}
% \Pic is the operator that assigns to X its picard group, usage \Pic(X)
% \Picardstack_{X/B} denotes the Picard stack of X over B
% \Picardfunctor_{X/B} denotes the Picard functor of X over B
\def\Pic{\mathop{\mathrm{Pic}}\nolimits}
\def\Picardstack{\mathcal{P}\!\mathit{ic}}
\def\Picardfunctor{\mathrm{Pic}}
\def\Deformationcategory{\mathcal{D}\!\mathit{ef}}
\newenvironment{reference}{}{}
\newenvironment{slogan}{}{}
\newenvironment{history}{}{}
\usepackage{cleveref}
\crefname{originallemma}{original lemma}{original lemmas}
\Crefname{originallemma}{Original lemma}{Original lemmas}
\makeatletter
\newlabel{lemma-chow-limit}{{0FH6}{source}{}{sourcepointer-1}{}}
\newlabel{lemma-chow-limit@cref}{{[originallemma][0][]0FH6}{[1][0][]source}{}{}{}}
\makeatother
\crefname{originallemma}{original lemma}{original lemmas}
\Crefname{originallemma}{Original lemma}{Original lemmas}
\makeatletter
\newlabel{varieties-lemma-generate-over-complement}{{0FD5}{source}{}{sourcepointer-2}{}}
\newlabel{varieties-lemma-generate-over-complement@cref}{{[originallemma][0][]0FD5}{[1][0][]source}{}{}{}}
\makeatother
\crefname{originallemma}{original lemma}{original lemmas}
\Crefname{originallemma}{Original lemma}{Original lemmas}
\makeatletter
\newlabel{varieties-lemma-bertini}{{0FD6}{source}{}{sourcepointer-3}{}}
\newlabel{varieties-lemma-bertini@cref}{{[originallemma][0][]0FD6}{[1][0][]source}{}{}{}}
\makeatother
\crefname{originallemma}{original lemma}{original lemmas}
\Crefname{originallemma}{Original lemma}{Original lemmas}
\makeatletter
\newlabel{varieties-lemma-connectedness-ample-divisor}{{0FD9}{source}{}{sourcepointer-4}{}}
\newlabel{varieties-lemma-connectedness-ample-divisor@cref}{{[originallemma][0][]0FD9}{[1][0][]source}{}{}{}}
\makeatother
\crefname{originallemma}{original lemma}{original lemmas}
\Crefname{originallemma}{Original lemma}{Original lemmas}
\makeatletter
\newlabel{pic-lemma-picard-pieces}{{0BA0}{source}{}{sourcepointer-5}{}}
\newlabel{pic-lemma-picard-pieces@cref}{{[originallemma][0][]0BA0}{[1][0][]source}{}{}{}}
\makeatother
\crefname{originallemma}{original lemma}{original lemmas}
\Crefname{originallemma}{Original lemma}{Original lemmas}
\makeatletter
\newlabel{varieties-lemma-geometrically-connected-criterion}{{056R}{source}{}{sourcepointer-6}{}}
\newlabel{varieties-lemma-geometrically-connected-criterion@cref}{{[originallemma][0][]056R}{[1][0][]source}{}{}{}}
\makeatother
\crefname{originallemma}{original lemma}{original lemmas}
\Crefname{originallemma}{Original lemma}{Original lemmas}
\makeatletter
\newlabel{varieties-lemma-smooth-geometrically-normal}{{056T}{source}{}{sourcepointer-7}{}}
\newlabel{varieties-lemma-smooth-geometrically-normal@cref}{{[originallemma][0][]056T}{[1][0][]source}{}{}{}}
\makeatother
\crefname{originaldefinition}{original definition}{original definitions}
\Crefname{originaldefinition}{Original definition}{Original definitions}
\makeatletter
\newlabel{groupoids-definition-abelian-variety}{{03RO}{source}{}{sourcepointer-8}{}}
\newlabel{groupoids-definition-abelian-variety@cref}{{[originaldefinition][0][]03RO}{[1][0][]source}{}{}{}}
\makeatother
\crefname{originalproposition}{original proposition}{original propositions}
\Crefname{originalproposition}{Original proposition}{Original propositions}
\makeatletter
\newlabel{groupoids-proposition-review-abelian-varieties}{{03RP}{source}{}{sourcepointer-9}{}}
\newlabel{groupoids-proposition-review-abelian-varieties@cref}{{[originalproposition][0][]03RP}{[1][0][]source}{}{}{}}
\makeatother
\crefname{originallemma}{original lemma}{original lemmas}
\Crefname{originallemma}{Original lemma}{Original lemmas}
\makeatletter
\newlabel{chow-lemma-finite-flat}{{02RH}{source}{}{sourcepointer-10}{}}
\newlabel{chow-lemma-finite-flat@cref}{{[originallemma][0][]02RH}{[1][0][]source}{}{}{}}
\makeatother
\crefname{originalsection}{original section}{original sections}
\Crefname{originalsection}{Original section}{Original sections}
\makeatletter
\newlabel{pic-section-hilbert-scheme-points}{{0B94}{source}{}{sourcepointer-11}{}}
\newlabel{pic-section-hilbert-scheme-points@cref}{{[originalsection][0][]0B94}{[1][0][]source}{}{}{}}
\makeatother
\crefname{originalsection}{original section}{original sections}
\Crefname{originalsection}{Original section}{Original sections}
\makeatletter
\newlabel{pic-section-divisors}{{0B9C}{source}{}{sourcepointer-12}{}}
\newlabel{pic-section-divisors@cref}{{[originalsection][0][]0B9C}{[1][0][]source}{}{}{}}
\makeatother
\crefname{originallemma}{original lemma}{original lemmas}
\Crefname{originallemma}{Original lemma}{Original lemmas}
\makeatletter
\newlabel{pic-lemma-universal-object}{{0B9G}{source}{}{sourcepointer-13}{}}
\newlabel{pic-lemma-universal-object@cref}{{[originallemma][0][]0B9G}{[1][0][]source}{}{}{}}
\makeatother
\crefname{originallemma}{original lemma}{original lemmas}
\Crefname{originallemma}{Original lemma}{Original lemmas}
\makeatletter
\newlabel{divisors-lemma-finite-locally-free-has-norm}{{0BD2}{source}{}{sourcepointer-14}{}}
\newlabel{divisors-lemma-finite-locally-free-has-norm@cref}{{[originallemma][0][]0BD2}{[1][0][]source}{}{}{}}
\makeatother
\crefname{originallemma}{original lemma}{original lemmas}
\Crefname{originallemma}{Original lemma}{Original lemmas}
\makeatletter
\newlabel{divisors-lemma-norm-in-normal-case}{{0BD3}{source}{}{sourcepointer-15}{}}
\newlabel{divisors-lemma-norm-in-normal-case@cref}{{[originallemma][0][]0BD3}{[1][0][]source}{}{}{}}
\makeatother
\crefname{originallemma}{original lemma}{original lemmas}
\Crefname{originallemma}{Original lemma}{Original lemmas}
\makeatletter
\newlabel{divisors-lemma-norm-invertible}{{0BCY}{source}{}{sourcepointer-16}{}}
\newlabel{divisors-lemma-norm-invertible@cref}{{[originallemma][0][]0BCY}{[1][0][]source}{}{}{}}
\makeatother
\crefname{originallemma}{original lemma}{original lemmas}
\Crefname{originallemma}{Original lemma}{Original lemmas}
\makeatletter
\newlabel{chow-lemma-vanish-above-dimension}{{0FDX}{source}{}{sourcepointer-17}{}}
\newlabel{chow-lemma-vanish-above-dimension@cref}{{[originallemma][0][]0FDX}{[1][0][]source}{}{}{}}
\makeatother
\crefname{originallemma}{original lemma}{original lemmas}
\Crefname{originallemma}{Original lemma}{Original lemmas}
\makeatletter
\newlabel{chow-lemma-chow-ring-projective-bundle}{{02TX}{source}{}{sourcepointer-18}{}}
\newlabel{chow-lemma-chow-ring-projective-bundle@cref}{{[originallemma][0][]02TX}{[1][0][]source}{}{}{}}
\makeatother
\crefname{originallemma}{original lemma}{original lemmas}
\Crefname{originallemma}{Original lemma}{Original lemmas}
\makeatletter
\newlabel{divisors-lemma-Pic-projective-space-UFD}{{0BXJ}{source}{}{sourcepointer-19}{}}
\newlabel{divisors-lemma-Pic-projective-space-UFD@cref}{{[originallemma][0][]0BXJ}{[1][0][]source}{}{}{}}
\makeatother
\crefname{originallemma}{original lemma}{original lemmas}
\Crefname{originallemma}{Original lemma}{Original lemmas}
\makeatletter
\newlabel{perfect-lemma-chi-locally-constant-geometric}{{0B9T}{source}{}{sourcepointer-20}{}}
\newlabel{perfect-lemma-chi-locally-constant-geometric@cref}{{[originallemma][0][]0B9T}{[1][0][]source}{}{}{}}
\makeatother
\crefname{originallemma}{original lemma}{original lemmas}
\Crefname{originallemma}{Original lemma}{Original lemmas}
\makeatletter
\newlabel{coherent-lemma-cohomology-projective-space-over-ring}{{01XT}{source}{}{sourcepointer-21}{}}
\newlabel{coherent-lemma-cohomology-projective-space-over-ring@cref}{{[originallemma][0][]01XT}{[1][0][]source}{}{}{}}
\makeatother
\title{Smash nilpotence up to torsion of differences of rational points}
\author{The Stacks Project Authors}
\date{stacks-project:weil.tex:local-snapshot-20261011; GNU FDL 1.2 or later}
\begin{document}
\maketitle
\noindent\textit{Editorial note.} The displayed expansion of (sum c\_1(L\_i))\textasciicircum{}n has coefficient n!, not the printed n\textasciicircum{}n. Thus the displayed last coefficient should be ((n-1)!)\textasciicircum{}n n!, still a nonzero integer. The claimed torsion conclusion uses only a nonzero multiple and is unchanged. The original displays are retained.\par\medskip
\noindent\textit{Editorial note.} In the family paragraph, the symbols M\_p and M\_q refer to the constructed L-prime fibres. This is a visible notation slip, not an additional line bundle.\par\medskip
\noindent\textit{Editorial note.} The norm identity and projective-bundle description have author-marked omitted routine details. They are covered by the explicitly named standard norm and Abel–Jacobi inputs; no alternate source proof is inserted.\par\medskip
\noindent\textit{Allowed standard prerequisites (editorial).}\begin{itemize}
\item Explicit prerequisite boundary. Bertini through finitely many smooth points, Chow pullback/pushforward, norms of line bundles, symmetric powers and Abel–Jacobi projective bundles for n>=2g-1, projective-bundle Picard formula, and Euler characteristic in connected proper flat families.
\end{itemize}\medskip
\ifdefined\hypertarget\hypertarget{sourcepointer-1}{}\fi
\noindent\textit{Source reference.} Original lemma 0FH6: focused-seeds/AGP/sources/stacks-snapshot-20261011/tags.txt, source lines 19346--19346. Stable Stacks tag for an expressly allowed foundational or notational input; the counted construction remains the two concrete core obligations. The retained local core is not replaced by this pointer. See standard\_prerequisites.\par\medskip
\ifdefined\hypertarget\hypertarget{sourcepointer-2}{}\fi
\noindent\textit{Source reference.} Original lemma 0FD5: focused-seeds/AGP/sources/stacks-snapshot-20261011/tags.txt, source lines 19205--19205. Stable Stacks tag for an expressly allowed foundational or notational input; the counted construction remains the two concrete core obligations. The retained local core is not replaced by this pointer. See standard\_prerequisites.\par\medskip
\ifdefined\hypertarget\hypertarget{sourcepointer-3}{}\fi
\noindent\textit{Source reference.} Original lemma 0FD6: focused-seeds/AGP/sources/stacks-snapshot-20261011/tags.txt, source lines 19206--19206. Stable Stacks tag for an expressly allowed foundational or notational input; the counted construction remains the two concrete core obligations. The retained local core is not replaced by this pointer. See standard\_prerequisites.\par\medskip
\ifdefined\hypertarget\hypertarget{sourcepointer-4}{}\fi
\noindent\textit{Source reference.} Original lemma 0FD9: focused-seeds/AGP/sources/stacks-snapshot-20261011/tags.txt, source lines 19209--19209. Stable Stacks tag for an expressly allowed foundational or notational input; the counted construction remains the two concrete core obligations. The retained local core is not replaced by this pointer. See standard\_prerequisites.\par\medskip
\ifdefined\hypertarget\hypertarget{sourcepointer-5}{}\fi
\noindent\textit{Source reference.} Original lemma 0BA0: focused-seeds/AGP/sources/stacks-snapshot-20261011/tags.txt, source lines 14193--14193. Stable Stacks tag for an expressly allowed foundational or notational input; the counted construction remains the two concrete core obligations. The retained local core is not replaced by this pointer. See standard\_prerequisites.\par\medskip
\ifdefined\hypertarget\hypertarget{sourcepointer-6}{}\fi
\noindent\textit{Source reference.} Original lemma 056R: focused-seeds/AGP/sources/stacks-snapshot-20261011/tags.txt, source lines 6719--6719. Stable Stacks tag for an expressly allowed foundational or notational input; the counted construction remains the two concrete core obligations. The retained local core is not replaced by this pointer. See standard\_prerequisites.\par\medskip
\ifdefined\hypertarget\hypertarget{sourcepointer-7}{}\fi
\noindent\textit{Source reference.} Original lemma 056T: focused-seeds/AGP/sources/stacks-snapshot-20261011/tags.txt, source lines 6721--6721. Stable Stacks tag for an expressly allowed foundational or notational input; the counted construction remains the two concrete core obligations. The retained local core is not replaced by this pointer. See standard\_prerequisites.\par\medskip
\ifdefined\hypertarget\hypertarget{sourcepointer-8}{}\fi
\noindent\textit{Source reference.} Original definition 03RO: focused-seeds/AGP/sources/stacks-snapshot-20261011/tags.txt, source lines 4937--4937. Stable Stacks tag for an expressly allowed foundational or notational input; the counted construction remains the two concrete core obligations. The retained local core is not replaced by this pointer. See standard\_prerequisites.\par\medskip
\ifdefined\hypertarget\hypertarget{sourcepointer-9}{}\fi
\noindent\textit{Source reference.} Original proposition 03RP: focused-seeds/AGP/sources/stacks-snapshot-20261011/tags.txt, source lines 4938--4938. Stable Stacks tag for an expressly allowed foundational or notational input; the counted construction remains the two concrete core obligations. The retained local core is not replaced by this pointer. See standard\_prerequisites.\par\medskip
\ifdefined\hypertarget\hypertarget{sourcepointer-10}{}\fi
\noindent\textit{Source reference.} Original lemma 02RH: focused-seeds/AGP/sources/stacks-snapshot-20261011/tags.txt, source lines 3624--3624. Stable Stacks tag for an expressly allowed foundational or notational input; the counted construction remains the two concrete core obligations. The retained local core is not replaced by this pointer. See standard\_prerequisites.\par\medskip
\ifdefined\hypertarget\hypertarget{sourcepointer-11}{}\fi
\noindent\textit{Source reference.} Original section 0B94: focused-seeds/AGP/sources/stacks-snapshot-20261011/tags.txt, source lines 14162--14162. Stable Stacks tag for an expressly allowed foundational or notational input; the counted construction remains the two concrete core obligations. The retained local core is not replaced by this pointer. See standard\_prerequisites.\par\medskip
\ifdefined\hypertarget\hypertarget{sourcepointer-12}{}\fi
\noindent\textit{Source reference.} Original section 0B9C: focused-seeds/AGP/sources/stacks-snapshot-20261011/tags.txt, source lines 14170--14170. Stable Stacks tag for an expressly allowed foundational or notational input; the counted construction remains the two concrete core obligations. The retained local core is not replaced by this pointer. See standard\_prerequisites.\par\medskip
\ifdefined\hypertarget\hypertarget{sourcepointer-13}{}\fi
\noindent\textit{Source reference.} Original lemma 0B9G: focused-seeds/AGP/sources/stacks-snapshot-20261011/tags.txt, source lines 14174--14174. Stable Stacks tag for an expressly allowed foundational or notational input; the counted construction remains the two concrete core obligations. The retained local core is not replaced by this pointer. See standard\_prerequisites.\par\medskip
\ifdefined\hypertarget\hypertarget{sourcepointer-14}{}\fi
\noindent\textit{Source reference.} Original lemma 0BD2: focused-seeds/AGP/sources/stacks-snapshot-20261011/tags.txt, source lines 14300--14300. Stable Stacks tag for an expressly allowed foundational or notational input; the counted construction remains the two concrete core obligations. The retained local core is not replaced by this pointer. See standard\_prerequisites.\par\medskip
\ifdefined\hypertarget\hypertarget{sourcepointer-15}{}\fi
\noindent\textit{Source reference.} Original lemma 0BD3: focused-seeds/AGP/sources/stacks-snapshot-20261011/tags.txt, source lines 14301--14301. Stable Stacks tag for an expressly allowed foundational or notational input; the counted construction remains the two concrete core obligations. The retained local core is not replaced by this pointer. See standard\_prerequisites.\par\medskip
\ifdefined\hypertarget\hypertarget{sourcepointer-16}{}\fi
\noindent\textit{Source reference.} Original lemma 0BCY: focused-seeds/AGP/sources/stacks-snapshot-20261011/tags.txt, source lines 14296--14296. Stable Stacks tag for an expressly allowed foundational or notational input; the counted construction remains the two concrete core obligations. The retained local core is not replaced by this pointer. See standard\_prerequisites.\par\medskip
\ifdefined\hypertarget\hypertarget{sourcepointer-17}{}\fi
\noindent\textit{Source reference.} Original lemma 0FDX: focused-seeds/AGP/sources/stacks-snapshot-20261011/tags.txt, source lines 19232--19232. Stable Stacks tag for an expressly allowed foundational or notational input; the counted construction remains the two concrete core obligations. The retained local core is not replaced by this pointer. See standard\_prerequisites.\par\medskip
\ifdefined\hypertarget\hypertarget{sourcepointer-18}{}\fi
\noindent\textit{Source reference.} Original lemma 02TX: focused-seeds/AGP/sources/stacks-snapshot-20261011/tags.txt, source lines 3712--3712. Stable Stacks tag for an expressly allowed foundational or notational input; the counted construction remains the two concrete core obligations. The retained local core is not replaced by this pointer. See standard\_prerequisites.\par\medskip
\ifdefined\hypertarget\hypertarget{sourcepointer-19}{}\fi
\noindent\textit{Source reference.} Original lemma 0BXJ: focused-seeds/AGP/sources/stacks-snapshot-20261011/tags.txt, source lines 14982--14982. Stable Stacks tag for an expressly allowed foundational or notational input; the counted construction remains the two concrete core obligations. The retained local core is not replaced by this pointer. See standard\_prerequisites.\par\medskip
\ifdefined\hypertarget\hypertarget{sourcepointer-20}{}\fi
\noindent\textit{Source reference.} Original lemma 0B9T: focused-seeds/AGP/sources/stacks-snapshot-20261011/tags.txt, source lines 14186--14186. Stable Stacks tag for an expressly allowed foundational or notational input; the counted construction remains the two concrete core obligations. The retained local core is not replaced by this pointer. See standard\_prerequisites.\par\medskip
\ifdefined\hypertarget\hypertarget{sourcepointer-21}{}\fi
\noindent\textit{Source reference.} Original lemma 01XT: focused-seeds/AGP/sources/stacks-snapshot-20261011/tags.txt, source lines 2532--2532. Stable Stacks tag for an expressly allowed foundational or notational input; the counted construction remains the two concrete core obligations. The retained local core is not replaced by this pointer. See standard\_prerequisites.\par\medskip
\begin{lemma}[Voevodsky]
\label{lemma-smash-nilpotence}
\begin{reference}
\cite{nilpotence}
\end{reference}
Let $k$ be a field. Let $X$ be a geometrically irreducible
smooth projective scheme over $k$. Let $x, x' \in X$ be $k$-rational points.
For $n$ large enough the class of the zero cycle
$$
([x] - [x']) \times \ldots \times ([x] - [x']) \in
\CH_0(X^n)
$$
is torsion.
\end{lemma}

\begin{lemma}
\label{lemma-divide-difference-points}
Let $k$ be a field. Let $X$ be a geometrically irreducible
smooth projective scheme over $k$. Let $x, x' \in X$ be $k$-rational points.
Let $n$ be an integer invertible in $k$.
Then there exists a finite separable extension $k'/k$ such that
the pullback of $[x] - [x']$ to $X_{k'}$
is divisible by $n$ in $\CH_0(X_{k'})$.
\end{lemma}

\begin{proof}
Let $k'$ be a separable algebraic closure of $k$. Suppose that we can show
the pullback of $[x] - [x']$ to $X_{k'}$ is divisible by $n$ in
$\CH_0(X_{k'})$. Then we conclude by Lemma \ref{lemma-chow-limit}.
Thus we may and do assume $k$ is separably algebraically closed.

\medskip\noindent
Suppose $\dim(X) > 1$. Let $\mathcal{L}$ be an ample invertible sheaf on $X$.
Set
$$
V = \{s \in H^0(X, \mathcal{L}) \mid s(x) = 0\text{ and }s(x') = 0 \}
$$
After replacing $\mathcal{L}$ by a power we see that for
a general $v \in V$ the corresponding divisor $H_v \subset X$ is smooth
away from $x$ and $x'$, see
Varieties, Lemmas \ref{varieties-lemma-generate-over-complement} and
\ref{varieties-lemma-bertini}. To find $v$ we use that $k$ is infinite (being
separably algebraically closed).
If we choose $s$ general, then the image of $s$ in
$\mathfrak m_x\mathcal{L}_x/\mathfrak m_x^2\mathcal{L}_x$
will be nonzero, which implies that $H_v$ is smooth at $x$
(details omitted). Similarly for $x'$. Thus $H_v$ is smooth.
By Varieties, Lemma \ref{varieties-lemma-connectedness-ample-divisor}
(applied to the base change of everything
to the algebraic closure of $k$)
we see that $H_v$ is geometrically connected.
It suffices to prove the result for
$[x] - [x']$ seen as an element of $\CH_0(H_v)$.
In this way we reduce to the case of a curve.

\medskip\noindent
Assume $X$ is a curve. Then we see that $\mathcal{O}_X(x - x')$
defines a $k$-rational point $g$ of $J = \underline{\Pic}^0_{X/k}$, see
Picard Schemes of Curves, Lemma \ref{pic-lemma-picard-pieces}.
Recall that $J$ is a proper smooth variety over $k$
which is also a group scheme over $k$ (same reference).
Hence $J$ is geometrically integral
(see Varieties, Lemma \ref{varieties-lemma-geometrically-connected-criterion}
and \ref{varieties-lemma-smooth-geometrically-normal}).
In other words, $J$ is an abelian variety, see
Groupoids, Definition \ref{groupoids-definition-abelian-variety}.
Thus $[n] : J \to J$ is finite \'etale by
Groupoids, Proposition \ref{groupoids-proposition-review-abelian-varieties}
(this is where we use $n$ is invertible in $k$).
Since $k$ is separably closed we conclude that $g = [n](g')$
for some $g' \in J(k)$. If $\mathcal{L}$ is the degree $0$
invertible module on $X$ corresponding to $g'$, then we conclude
that $\mathcal{O}_X(x - x') \cong \mathcal{L}^{\otimes n}$ as desired.
\end{proof}

\begin{lemma}
\label{lemma-kernel-to-closure}
Let $K/k$ be an algebraic extension of fields.
Let $X$ be a finite type scheme over $k$.
The kernel of the map $\CH_i(X) \to \CH_i(X_K)$
constructed in Lemma \ref{lemma-chow-limit}
is torsion.
\end{lemma}

\begin{proof}
It clearly suffices to show that the kernel
of flat pullback $\CH_i(X) \to \CH_i(X_{k'})$
by $\pi : X_{k'} \to X$ is torsion
for any finite extension $k'/k$. This is clear because
$\pi_* \pi^* \alpha = [k' : k] \alpha$ by
Chow Homology, Lemma \ref{chow-lemma-finite-flat}.
\end{proof}

\begin{proof}
If we can show this after base change to the algebraic closure of $k$,
then the result follows over $k$ because the kernel of pullback
is torsion by Lemma \ref{lemma-kernel-to-closure}.
Hence we may and do assume $k$ is algebraically closed.

\medskip\noindent
Using Bertini we can choose a smooth curve $C \subset X$ passing through
$x$ and $x'$. See proof of Lemma \ref{lemma-divide-difference-points}.
Hence we may assume $X$ is a curve.

\medskip\noindent
Assume $X$ is a curve and $k$ is algebraically closed.
Write $S^n(X) = \underline{\Hilbfunctor}^n_{X/k}$ with notation as in
Picard Schemes of Curves, Sections \ref{pic-section-hilbert-scheme-points}
and \ref{pic-section-divisors}. There is a canonical morphism
$$
\pi : X^n \longrightarrow S^n(X)
$$
which sends the $k$-rational point $(x_1, \ldots, x_n)$ to the $k$-rational
point corresponding to the divisor $[x_1] + \ldots + [x_n]$ on $X$.
There is a faithful action of the symmetric group $S_n$ on $X^n$.
The morphism $\pi$ is $S_n$-invariant and the fibres of $\pi$ are
$S_n$-orbits (set theoretically). Finally, $\pi$ is finite flat of
degree $n!$, see Picard Schemes of Curves, Lemma
\ref{pic-lemma-universal-object}.

\medskip\noindent
Let $\alpha_n$ be the zero cycle on $X^n$ given by the formula in the
statement of the lemma. Let $\mathcal{L} = \mathcal{O}_X(x - x')$. Then
$c_1(\mathcal{L}) \cap [X] = [x] - [x']$. Thus
$$
\alpha_n = c_1(\mathcal{L}_1) \cap \ldots \cap c_1(\mathcal{L}_n) \cap [X^n]
$$
where $\mathcal{L}_i = \text{pr}_i^*\mathcal{L}$ and $\text{pr}_i : X^n \to X$
is the $i$th projection. By either
Divisors, Lemma \ref{divisors-lemma-finite-locally-free-has-norm} or
Divisors, Lemma \ref{divisors-lemma-norm-in-normal-case}
there is a norm for $\pi$. Set
$\mathcal{N} = \text{Norm}_\pi(\mathcal{L}_1)$, 
see Divisors, Lemma \ref{divisors-lemma-norm-invertible}. We have
$$
\pi^*\mathcal{N} =
(\mathcal{L}_1 \otimes \ldots \otimes \mathcal{L}_n)^{\otimes (n - 1)!}
$$
in $\Pic(X^n)$ by a calculation. Deails omitted; hint: this follows from
the fact that
$\text{Norm}_\pi : \pi_*\mathcal{O}_{X^n} \to \mathcal{O}_{S^n(X)}$
composed with the natural map $\pi_*\mathcal{O}_{S^n(X)} \to \mathcal{O}_{X^n}$
is equal to the product over all $\sigma \in S_n$ of the action of
$\sigma$ on $\pi_*\mathcal{O}_{X^n}$. Consider
$$
\beta_n = c_1(\mathcal{N})^n \cap [S^n(X)]
$$
in $\CH_0(S^n(X))$. Observe that
$c_1(\mathcal{L}_i) \cap c_1(\mathcal{L}_i) = 0$
because $\mathcal{L}_i$ is pulled back from a curve, see
Chow Homology, Lemma \ref{chow-lemma-vanish-above-dimension}. Thus we see that
\begin{align*}
\pi^*\beta_n
& =
((n - 1)!)^n
(\sum\nolimits_{i = 1, \ldots, n} c_1(\mathcal{L}_i))^n \cap [X^n] \\
& =
((n - 1)!)^n n^n 
c_1(\mathcal{L}_1) \cap \ldots \cap c_1(\mathcal{L}_n) \cap [X^n] \\
& =
(n!)^n \alpha_n
\end{align*}
Thus it suffices to show that $\beta_n$ is torsion.

\medskip\noindent
There is a canonical morphism
$$
f : S^n(X) \longrightarrow \underline{\Picardfunctor}^n_{X/k}
$$
See Picard Schemes of Curves, Lemma \ref{pic-lemma-picard-pieces}.
For $n \geq 2g - 1$ this morphism is a projective space bundle
(details omitted; compare with the
proof of Picard Schemes of Curves, Lemma \ref{pic-lemma-picard-pieces}).
The invertible sheaf $\mathcal{N}$ is trivial on the
fibres of $f$, see below. Thus by the projective space bundle formula
(Chow Homology, Lemma \ref{chow-lemma-chow-ring-projective-bundle})
we see that $\mathcal{N} = f^*\mathcal{M}$ for some invertible
module $\mathcal{M}$ on $\underline{\Picardfunctor}^n_{X/k}$.
Of course, then we see that
$$
c_1(\mathcal{N})^n = f^*(c_1(\mathcal{M})^n)
$$
is zero because $n > g = \dim(\underline{\Picardfunctor}^n_{X/k})$
and we can use Chow Homology, Lemma \ref{chow-lemma-vanish-above-dimension}
as before.

\medskip\noindent
We still have to show that $\mathcal{N}$ is trivial on a fibre $F$
of $f$. Since the fibres of $f$ are projective spaces and since
$\Pic(\mathbf{P}^m_k) = \mathbf{Z}$
(Divisors, Lemma \ref{divisors-lemma-Pic-projective-space-UFD}),
this can be shown by computing the degree of $\mathcal{N}$
on a line contained in the fibre. Instead we will prove it by
proving that $\mathcal{N}$ is algebraically
equivalent to zero. First we claim there is a connected finite type scheme $T$
over $k$, an invertible module $\mathcal{L}'$ on $T \times X$ and
$k$-rational points $p, q \in T$ such that
$\mathcal{M}_p \cong \mathcal{O}_X$ and $\mathcal{M}_q = \mathcal{L}$.
Namely, since $\mathcal{L} = \mathcal{O}_X(x - x')$ we can take
$T = X$, $p = x'$, $q = x$, and
$\mathcal{L}' = \mathcal{O}_{X \times X}(\Delta)
\otimes \text{pr}_2^*\mathcal{O}_X(-x')$.
Then we let $\mathcal{L}'_i$ on
$T \times X^n$ for $i = 1, \ldots, n$
be the pullback of $\mathcal{L}'$ by
$\text{id}_T \times \text{pr}_i : T \times X^n \to T \times X$.
Finally, we let
$\mathcal{N}' = \text{Norm}_{\text{id}_T \times \pi}(\mathcal{L}'_1)$
on $T \times S^n(X)$.
By construction we have $\mathcal{N}'_p = \mathcal{O}_{S^n(X)}$
and $\mathcal{N}'_q = \mathcal{N}$.
We conclude that
$$
\mathcal{N}'|_{T \times F}
$$
is an invertible module on $T \times F \cong T \times \mathbf{P}^m_k$
whose fibre over $p$ is the trivial invertible module and whose fibre
over $q$ is $\mathcal{N}|_F$. Since the euler characteristic
of the trivial bundle is $1$ and since this euler characteristic
is locally constant in families (Derived Categories of Schemes,
Lemma \ref{perfect-lemma-chi-locally-constant-geometric})
we conclude $\chi(F, \mathcal{N}^{\otimes s}|_F) = 1$
for all $s \in \mathbf{Z}$. This can happen only if
$\mathcal{N}|_F \cong \mathcal{O}_F$ (see
Cohomology of Schemes, Lemma
\ref{coherent-lemma-cohomology-projective-space-over-ring})
and the proof is complete. Some details omitted.
\end{proof}

\bibliographystyle{plain}
\bibliography{my}
\end{document}
